\documentclass[11pt]{article}
\usepackage[T1]{fontenc}
\usepackage[utf8]{inputenc}
\usepackage{lmodern}
\usepackage{microtype}
\usepackage[margin=0.72in,headheight=15pt]{geometry}
\usepackage[table]{xcolor}
\usepackage{array,booktabs,longtable,tabularx}
\usepackage{amssymb}
\usepackage{enumitem}
\usepackage[most]{tcolorbox}
\usepackage{fancyhdr}
\usepackage{titlesec}
\usepackage{needspace}
\usepackage{ragged2e}
\usepackage{hyperref}
\usepackage{bookmark}
\usepackage{lastpage}

\definecolor{Navy}{HTML}{17324D}
\definecolor{Teal}{HTML}{157A7A}
\definecolor{CueGray}{HTML}{EEF2F4}
\definecolor{SoftBlue}{HTML}{E9F3F8}
\definecolor{SoftGreen}{HTML}{EAF5F0}
\definecolor{SoftAmber}{HTML}{FFF4D6}
\definecolor{RuleGray}{HTML}{B7C3CA}
\definecolor{TextGray}{HTML}{263238}

\hypersetup{colorlinks=true,linkcolor=Navy,urlcolor=Teal,citecolor=Navy,pdfborder={0 0 0}}
\color{TextGray}
\setlength{\parindent}{0pt}
\setlength{\parskip}{4pt}
\setlength{\emergencystretch}{3em}
\hfuzz=0.2pt
\hbadness=2000
\setlist{nosep,leftmargin=1.4em}
\setlength{\LTpre}{4pt}
\setlength{\LTpost}{6pt}
\renewcommand{\arraystretch}{1.2}

\titleformat{\section}{\Large\bfseries\color{Navy}\RaggedRight}{\thesection}{0.55em}{}
\titleformat{\subsection}{\normalsize\bfseries\color{Teal}\RaggedRight}{\thesubsection}{0.5em}{}
\titlespacing*{\section}{0pt}{1.3ex plus .3ex}{.7ex}
\titlespacing*{\subsection}{0pt}{1.1ex plus .2ex}{.5ex}

\pagestyle{fancy}
\fancyhf{}
\lhead{\small\color{Navy}\textbf{Cybersecurity Cornell Notes}}
\rhead{\small\color{Teal}\ChapterMark}
\cfoot{\small\thepage\ of \pageref*{LastPage}}
\renewcommand{\headrulewidth}{0.5pt}

\newtcolorbox{orientationbox}{enhanced,breakable,colback=SoftBlue,colframe=Teal,boxrule=.7pt,arc=2mm,left=2mm,right=2mm,top=1.5mm,bottom=1.5mm}
\newtcolorbox{topicsummary}[1][]{enhanced,breakable,colback=CueGray,colframe=RuleGray,title={Topic Summary},fonttitle=\bfseries\color{Navy},boxrule=.5pt,arc=1.5mm,left=2mm,right=2mm,top=1mm,bottom=1mm,#1}
\newtcolorbox{defensebox}{enhanced,breakable,colback=SoftGreen,colframe=Teal,title={Defensive Guidance},fonttitle=\bfseries,boxrule=.7pt,arc=2mm}
\newtcolorbox{historybox}{enhanced,breakable,colback=SoftAmber,colframe=orange!70!black,title={Historical Context},fonttitle=\bfseries,boxrule=.7pt,arc=2mm}
\newtcolorbox{synthesisbox}{enhanced,breakable,colback=Navy!5,colframe=Navy,title={Chapter Synthesis},fonttitle=\bfseries,boxrule=.8pt,arc=2mm}
\newtcolorbox{takeawaybox}{enhanced,colback=Teal!8,colframe=Teal,title={One-Sentence Takeaway},fonttitle=\bfseries,boxrule=.9pt,arc=2mm}

\newcommand{\CornellHeader}{%
\rowcolor{Navy}\color{white}\bfseries Cue / Recall Prompt & \color{white}\bfseries Notes / Analysis\\\hline}
\newcommand{\CornellRow}[2]{%
\cellcolor{CueGray}\RaggedRight\textbf{#1} & \RaggedRight #2\\\hline}

\newcommand{\ChapterMark}{Chapter 88}
\title{\textcolor{Navy}{Chapter 88: Cyber Attack Process}\\[2mm]\large Cornell Notes}
\author{}
\date{}
\hypersetup{pdftitle={Chapter 88: Cyber Attack Process - Cornell Notes},pdfauthor={},pdfsubject={Cybersecurity study notes},pdfkeywords={attack, cyber, process, and target}}
\begin{document}
\maketitle
\thispagestyle{fancy}
\vspace{-1.5em}
\begin{orientationbox}
\textbf{Chapter orientation.} This chapter examines cyber attack process within the broader domain of critical infrastructure security. Its central problem is how to reason about attack, cyber, process, and target while keeping security objectives, operating assumptions, evidence, and recovery connected. A practitioner should approach the material through physical consequences, safety, availability, environmental dependencies, remote connectivity, maintenance, and recovery. Useful prerequisites include basic risk, threat, control, and systems concepts.
\end{orientationbox}
\section*{Learning Objectives}
\begin{itemize}
\item Explain the security purpose and scope of Cyber Attack Process.
\item Identify the assets, actors, assumptions, and trust boundaries associated with attack, cyber, and process.
\item Distinguish Introduction from What Is A Cyber Attack and explain where their responsibilities intersect.
\item Analyze how failures in attack, cyber, and process can affect confidentiality, integrity, availability, privacy, or safety.
\item Evaluate relevant controls through physical consequences, safety, availability, environmental dependencies, remote connectivity, maintenance, and recovery.
\item Audit an implementation using asset and dependency maps, physical-access records, sensor telemetry, safety constraints, maintenance logs, and recovery exercises.
\item Prioritize remediation according to exploitability, consequence, control strength, and recovery readiness.
\item Verify that corrective actions change observable risk rather than documentation alone.
\end{itemize}
\section*{Cornell Cue-and-Notes}
\Needspace{8\baselineskip}
\subsection*{Introduction}
\begin{longtable}{>{\RaggedRight\arraybackslash}p{0.29\textwidth}|>{\RaggedRight\arraybackslash}p{0.65\textwidth}}
\CornellHeader
\CornellRow{What security problem does Introduction address?}{\textbf{Core idea.} Treat Introduction as a structured part of Cyber Attack Process, with particular attention to attack, cyber, malicious, and process. \textbf{Analytical lens.} This topic establishes a component of the chapter's argument; connect its definitions and mechanisms to a testable security objective. For attack, cyber, and malicious, model safety and physical consequences alongside conventional information-security effects. \textbf{Security reasoning.} Identify what must remain protected, who or what can alter the expected state, which assumptions cross a trust boundary, and how failure changes confidentiality, integrity, availability, privacy, or safety. \textbf{Application.} The practitioner should evaluate cyber and physical failure together, enforce safe states, and test operational recovery. \textbf{Evidence.} Seek asset and dependency maps, physical-access records, sensor telemetry, safety constraints, maintenance logs, and recovery exercises.}
\end{longtable}
\begin{topicsummary}
Introduction is best remembered as the connection between attack, cyber, malicious, and process and the chapter's larger control objective. A defensible review states the assumption, expected behavior, evidence, failure consequence, and required response.
\end{topicsummary}
\Needspace{8\baselineskip}
\subsection*{What Is A Cyber Attack}
\begin{longtable}{>{\RaggedRight\arraybackslash}p{0.29\textwidth}|>{\RaggedRight\arraybackslash}p{0.65\textwidth}}
\CornellHeader
\CornellRow{Which assumptions matter in What Is A Cyber Attack?}{\textbf{Core idea.} Treat What Is A Cyber Attack as a structured part of Cyber Attack Process, with particular attention to attack, cyber, process, and target. \textbf{Analytical lens.} This topic is threat-centered: separate actor capability, reachable attack surface, prerequisites, execution path, and consequence. For attack, cyber, and process, model safety and physical consequences alongside conventional information-security effects. \textbf{Security reasoning.} Identify what must remain protected, who or what can alter the expected state, which assumptions cross a trust boundary, and how failure changes confidentiality, integrity, availability, privacy, or safety. \textbf{Application.} The practitioner should evaluate cyber and physical failure together, enforce safe states, and test operational recovery. \textbf{Evidence.} Seek asset and dependency maps, physical-access records, sensor telemetry, safety constraints, maintenance logs, and recovery exercises.}
\end{longtable}
\begin{topicsummary}
What Is A Cyber Attack is best remembered as the connection between attack, cyber, process, and target and the chapter's larger control objective. A defensible review states the assumption, expected behavior, evidence, failure consequence, and required response.
\end{topicsummary}
\Needspace{8\baselineskip}
\subsection*{The Cyber Attack Process}
\begin{longtable}{>{\RaggedRight\arraybackslash}p{0.29\textwidth}|>{\RaggedRight\arraybackslash}p{0.65\textwidth}}
\CornellHeader
\CornellRow{How should a reviewer evaluate The Cyber Attack Process?}{\textbf{Core idea.} Treat The Cyber Attack Process as a structured part of Cyber Attack Process, with particular attention to attack, process, cyber, and users. \textbf{Analytical lens.} This topic is threat-centered: separate actor capability, reachable attack surface, prerequisites, execution path, and consequence. For attack, process, and cyber, model safety and physical consequences alongside conventional information-security effects. \textbf{Security reasoning.} Identify what must remain protected, who or what can alter the expected state, which assumptions cross a trust boundary, and how failure changes confidentiality, integrity, availability, privacy, or safety. \textbf{Application.} The practitioner should evaluate cyber and physical failure together, enforce safe states, and test operational recovery. \textbf{Evidence.} Seek asset and dependency maps, physical-access records, sensor telemetry, safety constraints, maintenance logs, and recovery exercises. Related subtopics include Today's Complex Environments.}
\end{longtable}
\begin{topicsummary}
The Cyber Attack Process is best remembered as the connection between attack, process, cyber, and users and the chapter's larger control objective. A defensible review states the assumption, expected behavior, evidence, failure consequence, and required response.
\end{topicsummary}
\Needspace{8\baselineskip}
\subsection*{Attack Difficulty And Complexity}
\begin{longtable}{>{\RaggedRight\arraybackslash}p{0.29\textwidth}|>{\RaggedRight\arraybackslash}p{0.65\textwidth}}
\CornellHeader
\CornellRow{Why can Attack Difficulty And Complexity fail in practice?}{\textbf{Core idea.} Treat Attack Difficulty And Complexity as a structured part of Cyber Attack Process, with particular attention to malicious, attack, cyber, and level. \textbf{Analytical lens.} This topic is threat-centered: separate actor capability, reachable attack surface, prerequisites, execution path, and consequence. For malicious, attack, and cyber, model safety and physical consequences alongside conventional information-security effects. \textbf{Security reasoning.} Identify what must remain protected, who or what can alter the expected state, which assumptions cross a trust boundary, and how failure changes confidentiality, integrity, availability, privacy, or safety. \textbf{Application.} The practitioner should evaluate cyber and physical failure together, enforce safe states, and test operational recovery. \textbf{Evidence.} Seek asset and dependency maps, physical-access records, sensor telemetry, safety constraints, maintenance logs, and recovery exercises.}
\end{longtable}
\begin{topicsummary}
Attack Difficulty And Complexity is best remembered as the connection between malicious, attack, cyber, and level and the chapter's larger control objective. A defensible review states the assumption, expected behavior, evidence, failure consequence, and required response.
\end{topicsummary}
\Needspace{8\baselineskip}
\subsection*{Attack Vectors}
\begin{longtable}{>{\RaggedRight\arraybackslash}p{0.29\textwidth}|>{\RaggedRight\arraybackslash}p{0.65\textwidth}}
\CornellHeader
\CornellRow{What security problem does Attack Vectors address?}{\textbf{Core idea.} Treat Attack Vectors as a structured part of Cyber Attack Process, with particular attention to attack, target, vectors, and process. \textbf{Analytical lens.} This topic is threat-centered: separate actor capability, reachable attack surface, prerequisites, execution path, and consequence. For attack, target, and vectors, model safety and physical consequences alongside conventional information-security effects. \textbf{Security reasoning.} Identify what must remain protected, who or what can alter the expected state, which assumptions cross a trust boundary, and how failure changes confidentiality, integrity, availability, privacy, or safety. \textbf{Application.} The practitioner should evaluate cyber and physical failure together, enforce safe states, and test operational recovery. \textbf{Evidence.} Seek asset and dependency maps, physical-access records, sensor telemetry, safety constraints, maintenance logs, and recovery exercises.}
\end{longtable}
\begin{topicsummary}
Attack Vectors is best remembered as the connection between attack, target, vectors, and process and the chapter's larger control objective. A defensible review states the assumption, expected behavior, evidence, failure consequence, and required response.
\end{topicsummary}
\Needspace{8\baselineskip}
\subsection*{Beginning The Cyber Attack Process}
\begin{longtable}{>{\RaggedRight\arraybackslash}p{0.29\textwidth}|>{\RaggedRight\arraybackslash}p{0.65\textwidth}}
\CornellHeader
\CornellRow{Which assumptions matter in Beginning The Cyber Attack Process?}{\textbf{Core idea.} Treat Beginning The Cyber Attack Process as a structured part of Cyber Attack Process, with particular attention to access, attacker, target, and cyber. \textbf{Analytical lens.} This topic is threat-centered: separate actor capability, reachable attack surface, prerequisites, execution path, and consequence. For access, attacker, and target, model safety and physical consequences alongside conventional information-security effects. \textbf{Security reasoning.} Identify what must remain protected, who or what can alter the expected state, which assumptions cross a trust boundary, and how failure changes confidentiality, integrity, availability, privacy, or safety. \textbf{Application.} The practitioner should evaluate cyber and physical failure together, enforce safe states, and test operational recovery. \textbf{Evidence.} Seek asset and dependency maps, physical-access records, sensor telemetry, safety constraints, maintenance logs, and recovery exercises. Related subtopics include Access, Vulnerabilities, Exploits, and Malware.}
\end{longtable}
\begin{topicsummary}
Beginning The Cyber Attack Process is best remembered as the connection between access, attacker, target, and cyber and the chapter's larger control objective. A defensible review states the assumption, expected behavior, evidence, failure consequence, and required response.
\end{topicsummary}
\Needspace{8\baselineskip}
\subsection*{Cyber Attack Impacts}
\begin{longtable}{>{\RaggedRight\arraybackslash}p{0.29\textwidth}|>{\RaggedRight\arraybackslash}p{0.65\textwidth}}
\CornellHeader
\CornellRow{How should a reviewer evaluate Cyber Attack Impacts?}{\textbf{Core idea.} Treat Cyber Attack Impacts as a structured part of Cyber Attack Process, with particular attention to malware, attacks, software, and cyber. \textbf{Analytical lens.} This topic is threat-centered: separate actor capability, reachable attack surface, prerequisites, execution path, and consequence. For malware, attacks, and software, model safety and physical consequences alongside conventional information-security effects. \textbf{Security reasoning.} Identify what must remain protected, who or what can alter the expected state, which assumptions cross a trust boundary, and how failure changes confidentiality, integrity, availability, privacy, or safety. \textbf{Application.} The practitioner should evaluate cyber and physical failure together, enforce safe states, and test operational recovery. \textbf{Evidence.} Seek asset and dependency maps, physical-access records, sensor telemetry, safety constraints, maintenance logs, and recovery exercises. Related subtopics include Supply Chain Attack, and Target and Home Depot Reported in 2014.}
\end{longtable}
\begin{topicsummary}
Cyber Attack Impacts is best remembered as the connection between malware, attacks, software, and cyber and the chapter's larger control objective. A defensible review states the assumption, expected behavior, evidence, failure consequence, and required response.
\end{topicsummary}
\Needspace{8\baselineskip}
\subsection*{Defending Against Cyber Attacks}
\begin{longtable}{>{\RaggedRight\arraybackslash}p{0.29\textwidth}|>{\RaggedRight\arraybackslash}p{0.65\textwidth}}
\CornellHeader
\CornellRow{Why can Defending Against Cyber Attacks fail in practice?}{\textbf{Core idea.} Treat Defending Against Cyber Attacks as a structured part of Cyber Attack Process, with particular attention to cyber, attack, users, and devices. \textbf{Analytical lens.} This topic is threat-centered: separate actor capability, reachable attack surface, prerequisites, execution path, and consequence. For cyber, attack, and users, model safety and physical consequences alongside conventional information-security effects. \textbf{Security reasoning.} Identify what must remain protected, who or what can alter the expected state, which assumptions cross a trust boundary, and how failure changes confidentiality, integrity, availability, privacy, or safety. \textbf{Application.} The practitioner should evaluate cyber and physical failure together, enforce safe states, and test operational recovery. \textbf{Evidence.} Seek asset and dependency maps, physical-access records, sensor telemetry, safety constraints, maintenance logs, and recovery exercises. Related subtopics include Jeep Hack/Sprint Experiment, Layers, Check and Validate Everything, and Awareness.}
\end{longtable}
\begin{topicsummary}
Defending Against Cyber Attacks is best remembered as the connection between cyber, attack, users, and devices and the chapter's larger control objective. A defensible review states the assumption, expected behavior, evidence, failure consequence, and required response.
\end{topicsummary}
\begin{historybox}
Some products, protocols, examples, threat observations, or legal references in this chapter are historically bounded. Retain the underlying threat model and control objective, but verify present applicability before treating a named implementation as current guidance.
\end{historybox}
\begin{defensebox}
Apply the chapter by making assets, authority, trust, expected behavior, and failure response explicit. In practice, evaluate cyber and physical failure together, enforce safe states, and test operational recovery. A control is not fully supported until its operation, monitoring, ownership, and recovery behavior are demonstrated with evidence.
\end{defensebox}
\section*{Threat-and-Control Map}
\begingroup\scriptsize\setlength{\tabcolsep}{2pt}
\begin{longtable}{>{\RaggedRight\arraybackslash}p{0.135\textwidth}>{\RaggedRight\arraybackslash}p{0.145\textwidth}>{\RaggedRight\arraybackslash}p{0.155\textwidth}>{\RaggedRight\arraybackslash}p{0.145\textwidth}>{\RaggedRight\arraybackslash}p{0.16\textwidth}>{\RaggedRight\arraybackslash}p{0.17\textwidth}}
\toprule\rowcolor{Navy}\color{white}\textbf{Asset} & \color{white}\textbf{Threat / Failure} & \color{white}\textbf{Vulnerability} & \color{white}\textbf{Impact} & \color{white}\textbf{Detection / Evidence} & \color{white}\textbf{Control}\\\midrule
\endfirsthead
\toprule\rowcolor{Navy}\color{white}\textbf{Asset} & \color{white}\textbf{Threat / Failure} & \color{white}\textbf{Vulnerability} & \color{white}\textbf{Impact} & \color{white}\textbf{Detection / Evidence} & \color{white}\textbf{Control}\\\midrule
\endhead
People, physical processes, connected devices, and essential services & Tampering, unsafe command, service disruption, surveillance, or cascading failure & Insecure remote access, fragile dependencies, weak update paths, or insufficient safety separation & Safety events, prolonged outage, physical damage, privacy harm, or public disruption & Operational telemetry, integrity monitoring, physical controls, and anomaly investigation & Layered cyber-physical protection, safe-state design, segmentation, secure maintenance, and rehearsed recovery\\\midrule
Assumptions surrounding attack & Misuse or unexpected state transition & Trust in cyber is not validated & A local weakness becomes a broader compromise path & Review evidence for process and test negative paths & Make trust explicit, constrain authority, and fail safely\\\midrule
Recovery capability and decision evidence & Control failure remains unnoticed or cannot be reversed & Monitoring, ownership, or restoration has not been exercised & Longer exposure and slower, less reliable recovery & Alert-to-response exercises and restoration tests & Assigned ownership, actionable telemetry, preserved evidence, and rehearsed recovery\\\midrule
\bottomrule
\end{longtable}
\endgroup
\section*{Practical Security Checklist}
\begin{itemize}[label=$\square$]
\item Define the in-scope assets, users, services, data, and operational dependencies for Cyber Attack Process.
\item Document the expected behavior and security objective of attack.
\item Map identities, privileges, trust boundaries, and externally controlled inputs associated with attack and cyber.
\item Record the threat or failure being addressed, its prerequisites, likely path, and credible consequences.
\item Inspect asset and dependency maps, physical-access records, sensor telemetry, safety constraints, maintenance logs, and recovery exercises.
\item Test both the intended path and negative paths, including invalid state, partial failure, excessive privilege, and unavailable dependencies.
\item Confirm preventive controls are complemented by actionable detection and a named response owner.
\item Exercise containment, restoration, and process; record actual recovery behavior and unresolved dependencies.
\item Validate exceptions, compensating controls, expiration dates, and accountable approval.
\item Preserve reproducible evidence and tie each finding to affected assets, risk, remediation, and verification criteria.
\end{itemize}
\begin{synthesisbox}
The chapter's principal lesson is that cyber attack process must be evaluated as an operating security system rather than an isolated product or statement. The most important failure patterns involve attack, cyber, process, and target, especially when ownership, boundaries, telemetry, or recovery remain implicit. A reviewer should connect architecture and behavior to asset and dependency maps, physical-access records, sensor telemetry, safety constraints, maintenance logs, and recovery exercises, prioritize findings by reachable consequence and control independence, and verify remediation through observable change. Within Critical Infrastructure Security, this chapter contributes a reusable method: state the objective, expose assumptions, test failure, preserve evidence, and learn from operation.
\end{synthesisbox}
\section*{Self-Test Questions}
\begin{enumerate}
\item What is the central security objective of Cyber Attack Process?
\item How does Introduction contribute to the chapter's broader security argument?
\item Distinguish What Is A Cyber Attack from The Cyber Attack Process.
\item Which trust assumptions should be made explicit when evaluating attack?
\item How could a weakness involving cyber affect confidentiality, integrity, availability, privacy, or safety?
\item What evidence would demonstrate that controls surrounding process operate as intended?
\item Why should prevention, detection, response, and recovery be evaluated together in this domain?
\item Scenario: A connected physical process remains functionally available after a cyber event but no longer operates within safe constraints. What should the reviewer examine first, and why?
\item Scenario: A control associated with target passes a configuration check but fails during an operational exercise. How should the finding be interpreted and prioritized?
\item How should a practitioner distinguish enduring principles in this chapter from time-sensitive tools, examples, or implementation details?
\end{enumerate}
\Needspace{8\baselineskip}
\section*{Answer Key}
\begin{enumerate}
\item The objective is to protect the relevant assets by understanding physical consequences, safety, availability, environmental dependencies, remote connectivity, maintenance, and recovery and by connecting controls to measurable risk reduction.
\item Introduction provides one part of the causal chain: it should identify expected behavior, dependencies, failure modes, and the evidence needed to justify confidence.
\item What Is A Cyber Attack and The Cyber Attack Process address different portions of the problem. A strong comparison states their scope, assumptions, enforcement points, evidence, and how a failure in one affects the other.
\item Reviewers should identify who controls attack, what inputs or dependencies it trusts, where privilege changes, what happens during partial failure, and how those assumptions are tested.
\item A weakness in cyber can propagate across trust boundaries. The actual impact depends on reachable assets, granted authority, control independence, visibility, and recovery capability.
\item Useful evidence includes asset and dependency maps, physical-access records, sensor telemetry, safety constraints, maintenance logs, and recovery exercises; configuration alone is insufficient when operation, monitoring, or recovery is part of the claim.
\item Prevention reduces probability, detection limits dwell time, response limits spread, and recovery restores objectives. Omitting one layer creates an unexamined failure mode.
\item Begin by establishing scope, ownership, expected behavior, and the violated assumption. Then evaluate cyber and physical failure together, enforce safe states, and test operational recovery, preserving evidence for prioritization and follow-up.
\item Treat the exercise as stronger evidence than a static check. Prioritize according to reachable impact, likelihood of real operating conditions, missing compensating controls, and recovery difficulty.
\item Retain the principle, threat model, and control objective; separately label technologies, legal references, product examples, and threat observations whose applicability depends on time or environment.
\end{enumerate}
\section*{Key-Term Recap}
\begin{longtable}{>{\RaggedRight\arraybackslash}p{0.27\textwidth}>{\RaggedRight\arraybackslash}p{0.66\textwidth}}
\toprule\rowcolor{Navy}\color{white}\textbf{Term} & \color{white}\textbf{Meaning}\\\midrule
\endfirsthead
\toprule\rowcolor{Navy}\color{white}\textbf{Term} & \color{white}\textbf{Meaning}\\\midrule
\endhead
\textbf{Malware} & Software or code intentionally designed to disrupt, damage, spy, or obtain unauthorized control. \\\midrule
\textbf{Asset} & Information, service, capability, or physical resource whose value justifies protection. \\\midrule
\textbf{Authentication} & The process of establishing confidence that an asserted identity is genuine. \\\midrule
\textbf{Vulnerability} & A weakness that can be exploited or triggered to violate a security objective. \\\midrule
\textbf{Threat} & A circumstance or actor capable of causing harm by exploiting a weakness or adverse condition. \\\midrule
\textbf{Risk} & The effect of uncertainty on objectives, commonly evaluated through likelihood, impact, exposure, and context. \\\midrule
\textbf{Endpoint} & A user or server device that executes workloads and participates in a networked environment. \\\midrule
\textbf{Integrity} & The property that information and systems remain accurate, complete, and protected from unauthorized modification. \\\midrule
\textbf{Authorization} & The decision process that determines which authenticated subject may perform which action on a resource. \\\midrule
\textbf{Availability} & The property that authorized users can access information and services when required. \\\midrule
\bottomrule
\end{longtable}
\begin{takeawaybox}
Treat cyber attack process as a verifiable chain from assets and assumptions through controls, evidence, response, and recovery.
\end{takeawaybox}
\end{document}
